\documentclass{article}
\usepackage[UTF8]{ctex}
\usepackage{amsmath,amssymb,amsfonts}
\usepackage{enumitem}
\usepackage{geometry}
\geometry{a4paper, margin=1in}
% === 错题本排版优化（不参与题目标识） ===
\setlength{\parindent}{0pt}
\setlength{\parskip}{0.35em plus 0.10em minus 0.08em}
\setlength{\emergencystretch}{3em}
\setlist{itemsep=0.25em, topsep=0.40em, parsep=0pt}
\allowdisplaybreaks[2]
% === 错题本排版优化结束 ===
\begin{document}
\section*{广西大学2026年高等代数真题}
\section*{154．广西大学2026年高等代数试卷}
三。证明题。每题15分，共90分。\\

\subsection{题目1}
1．设 \(n\) 级方阵 \(A\) 满足 \(A^{2}=A\) 且 \(A \neq I\) ，问 \(A\) 是否可逆？请简述理由。角答题、每题15分，共30分。

\subsection{题目2}
2．请简述艾森斯坦因判别法，并且用其判断多项式 \(g(z)=z^{10}+z^{9}+\cdots+z+1\) 在有理数域中是否可以。\\
二计算题．每题15分，共30分．

\subsection{题目3}
3．计算下列 \(n\) 级行列式的值：

\[
\left|\begin{array}{cccc}
a+x_{1} & a & \cdots & a \\
a & a+x_{2} & \cdots & a \\
\vdots & \vdots & & \vdots \\
a & a & \cdots & a
\end{array}\right|, \text { 其中 } x_{1} x_{2} \cdots x_{n} \neq 0, n \geq 2 \text {. }
\]

\subsection{题目4}
4．设 \(\mathbb{R}^{4}\) 中的子空间 \(V_{1}\) 和 \(V_{2}\) 分别是下面两个线性方程组的解空间

\[
\left\{\begin{array} { l } 
{ x _ { 1 } - x _ { 2 } + x _ { 3 } = 0 ; } \\
{ 2 x _ { 1 } - x _ { 2 } + x _ { 3 } - x _ { 4 } = 0 . }
\end{array} \left\{\begin{array}{l}
x_{1}+x_{2}-x_{3}+x_{4}=0 ; \\
x_{2}-x_{3}+2 x_{4}=0 .
\end{array}\right.\right.
\]

\begin{itemize}
  \item[（1）] 求子空间 \(V_{1}+V_{2}\) 的维数和一个基。
  \item[（2）] 求子空间 \(V_{1} \cap V_{2}\) 的维数和一个基．
\end{itemize}

\subsection{题目5}
5．设 \(\mathbb{R}[x]_{n}\) 是实数域上次数小于等于 \(n\) 的多项式函数构成的线性空间，在 \(\mathbb{R}[x]_{n}\) 中定义求导变换 \(\mathcal{D}\)为

\[
\begin{aligned}
\mathcal{D}: \mathbb{R}[x]_{n} & \rightarrow \mathbb{R}[x]_{n} \\
a_{n} x^{n}+\cdots+a_{1} x+a_{0} & \mapsto n a_{n} x^{n-1}+\cdots+a_{1} .
\end{aligned}
\]

（1）证明： \(1, x, x^{2}, \cdots, x^{n}\) 是 \(\mathbb{R}[x]_{n}\) 的一个基，并写出 \(\mathcal{D}\) 在 \(1, x, x^{2}, \cdots, x^{n}\) 下的矩阵表示。\\
（2）证明： \(\mathcal{D}\) 是 \(\mathbb{R}[x]_{n}\) 上的幂零指数为 \(n+1\) 的幂零变换。

\subsection{题目6}
6．设 \(A\) 是 \(n\) 级实对称矩阵 \((n \geq 3), A^{*}\) 表示 \(A\) 的伴随矩阵，证明：\(\left(A^{*}\right)^{*}=|A|^{n-2} A\) ．

\subsection{题目7}
7．设 \(A\) 是 \(n\) 级实对称矩阵，证明：\(A\) 是半正定矩阵当且仅当对任意的正整数 \(m\) ，都存在唯一的实半正定矩阵 \(B\) ，使得 \(B^{m}=A\) 。

\subsection{题目8}
8．证明：在数域 \(K\) 上的一元多项式环 \(K[x]\) 中，\(d(x)\) 是 \(f(x)\) 和 \(g(x)\) 的最大公因式当且仅当存在\(u(x), v(x) \in K[x]\) 使得 \(d(x)=u(x) f(x)+v(x) g(x)\) 且 \(d(x)\) 是 \(f(x)\) 和 \(g(x)\) 的公因式．

\subsection{题目9}
9．设 \(A\) 是数域 \(K\) 上的 \(s \times n\) 矩阵，且 \(\operatorname{rank}(A)=n-t\) ，已知 \(\gamma_{1}, \gamma_{2}, \cdots, \gamma_{t+1}\) 是非齐次线性方程组\(A X=\alpha\) 的解，且非齐次线性方程组 \(A X=\alpha\) 的任意一个解 \(\gamma\) 都可以表示为

\[
\gamma=u_{1} \gamma_{1}+u_{2} \gamma_{2}+\cdots+u_{t+1} \gamma_{r+1} \text {, 其中 } u_{1}+u_{2}+\cdots+u_{t+1}=1 \text {. }
\]

证明：\(\gamma_{2}-\gamma_{1}, \gamma_{3}-\gamma_{1}, \cdots, \gamma_{t+1}-\gamma_{1}\) 是导出组 \(A X=0\) 的一个基础解系。

\subsection{题目10}
10．设 \(\alpha_{1}, \alpha_{2}, \cdots, \alpha_{m}\) 和 \(\beta_{1}, \beta_{2}, \cdots, \beta_{m}\) 是 \(n\) 维欧几里得空间的两个向量组，证明：存在一个正交变换\(\mathcal{A}\) 使得 \(\mathcal{A} \alpha_{i}=\beta_{i}(i=1,2, \cdots, m)\) 的充分必要条件为 \(\left(\alpha_{i}, \alpha_{j}\right)=\left(\beta_{i}, \beta_{j}\right)(i, j=1,2, \cdots, m)\) 。
\end{document}
